


\def\answerline{ % double horizontal line placed 0.5 cm below text, space between lines is 0.07 cm, then 0.75 cm of space 
	\vspace{0.5 cm}
	\hrule
	\vspace{0.07 cm}
	\hrule
	\vspace{0.75 cm}\noindent}
\NewDocumentCommand\length{O{3pt}}{\setlength\jot{#1}}% for align vertical spacing, there's probably a better way to do this locally


\NewDocumentCommand\thetitle{m O{}}{ \title{ \hypertarget{top}{\textbf{#1}} \\ \large {#2} } } %Bold title that is a hypertarget, optional bolded subtitle that's scaled properly 

% Transverse Quantities
\NewDocumentCommand{\pT}{}{$p_{\rm T}$} % transverse momentum
\NewDocumentCommand{\ET}{}{$E_{\rm T}$} % transverse Energy
\NewDocumentCommand{\MET}{}{$E_{\rm T}^{\rm miss}$} % transverse Energy

% Particles
\NewDocumentCommand{\h}{s}{\IfBooleanTF{#1}{\(H\)}{\(H^0\)}}
\NewDocumentCommand{\Z}{s}{\IfBooleanTF{#1}{\(Z\)}{\(Z^0\)}}
\NewDocumentCommand{\W}{s O{+}}{\IfBooleanTF{#1}{\(W^{#2}\)}{\(W^{\pm}\)}}

% Math symbols
\NewDocumentCommand{\e}{s m}{ % basis vector command
	\IfBooleanTF{#1} % \e{x} for x hat notation, \e*{x} for e_x notation
	{\mathbf{\hat{e}_{#2}}}
	{\bm{\hat{#2}}}} % from package bm, more powerful than \mathbf} 
\def\ints{\int_\mathcal{S}} % Surface integral
\def\intv{\int_\mathcal{V}} % Volume integral with V subscript
\def\intall{\int_{-\infty}^{\infty}} % Integral over all space
\def\iintall{\iint_{\rm All Space}} % Integral over all space
\def\iiintall{\iiint_{\rm All Space}} % Integral over all space
\def\ik{4\pi\epsilon_0} % Inverse k for EM
\def\lap{\mathcal{L}}% Laplace
\def\d{{\rm d}} % for integration
\NewDocumentCommand\slsh{m}{\slashed{#1}}
\NewDocumentCommand\dl{s}{\IfBooleanTF{#1}{}{\cdot} \d\mathbf{l}} % quicker curve integral dl, if no star, it also makes the dot 
\NewDocumentCommand\da{s}{\IfBooleanTF{#1}{}{\cdot} \d\mathbf{a}} % quicker surface integral da, if no star, it also makes the dot 
\NewDocumentCommand\oo{O{1} m} {\frac{#1}{#2}}   % reciprocal- 'one over', option to make it a regular fraction because oo is still quicker to type.

% Math Structures
\NewDocumentCommand{\parametric}{m}{ %command for getting boundary conditions with a "{" on the left
    \left\{ 
	\begin{gathered}
		\begin{matrix}
			#1
    	\end{matrix}
	\end{gathered}
	 \right.% nothing shows up on the right side
}% Don't trust the highlighting, this has the correct number of curly braces.

\NewDocumentCommand{\colvec}{+m}{
	\begingroup
		\renewcommand{\arraystretch}{1.25} % Default is 1.2. Adjust vertical spacing
		\begin{gathered}
			\begin{bmatrix}
				#1
			\end{bmatrix}
		\end{gathered}
	\endgroup
}
\NewDocumentCommand{\der}{s O{} m g}{% custom derivative command
	% \pder*{x} 		--> d_x
	% \pder[2]{x}		--> d^{2}f/d^{2}x
	% \pder{x}			--> d/dx
	% \pder{f}{x}		--> df/dx
	% \pder{f}{x}{y}	--> (df/dx)_y
	% \pder*[a]{b}{c}{e}--> (d^{a}_{c} b)_{e}
	% \pder*[a]{b}{c}{e}--> (d^{a}_{c} b)_{e}
     \IfBooleanTF{#1}
    {\IfNoValueTF{#4}  % g returns -NoValue- if no argument (read below)
        {\mathrm{D}_{#3}^{#2}} 
        {\mathrm{D}_{#4}^{#2}#3}
            }
    {\IfNoValueTF{#4}
        {\frac{\mathrm{d}^{#2}}{\mathrm{d} #3^{#2}}}
        {\frac{\mathrm{d}^{#2} #3}{\mathrm{d} #4^{#2}}}
	}}
\NewDocumentCommand{\pder}{s O{} m g d[]}{% custom partial derivative command, 
	% \pder*{x} 		--> ∂_x
	% \pder[2]{x}		--> ∂^{2}f/∂^{2}x
	% \pder{x}			--> ∂/∂x
	% \pder{f}{x}		--> ∂f/∂x
	% \pder{f}{x}{y}	--> (∂f/∂x)_y
	% \pder*[a]{b}{c}[e]--> (∂^{a}_{c} b)_{e}
	% \pder[a]{b}{c}[e]--> (∂^{a}b/∂^{a}c)_{e}
	\IfBooleanTF{#1}
	   {\IfNoValueTF{#4} 
		   {\partial_{#3}^{#2}}
		   {\partial_{#4}^{#2}#3}
			   }
	   {\IfNoValueTF{#5}% d[] returns -NoValue- if no argument present delimited by []
		   {\IfNoValueTF{#4}% g returns -NoValue- if no argument is present
			   {\frac{\partial^{#2}}{\partial #3^{#2}}}
			   {\frac{\partial^{#2} #3}{\partial #4^{#2}}}}
		   {{\IfNoValueTF{#4} 
			   {\left(\frac{\partial^{#2}}{\partial #3^{#2}}\right)_{\!#5}\!\!} % \! is thin negative horizontal space
			   {\left(\frac{\partial^{#2} #3}{\partial #4^{#2}}\right)_{\!#5}\!\!}}}
			   }} 
                
\NewDocumentCommand{\header}{m m m}{% Makes header that links to the key "top", usually used for first page
	\fancyhead[L]{#1} 
	\fancyhead[C]{\hyperlink{top}{ \textbf{#2} }}
	\fancyhead[R]{#3}
	\fancyfoot[C]{--\thepage--}
	\pagestyle{fancy}}
\setlength\headheight{17pt}
\addtolength{\topmargin}{-2.49998pt}

    
\NewDocumentCommand{\coloredanswer}{O{LavenderBlush2} m}{% makes coloredbox with pretty color
	\mathchoice
	{\colorbox{#1}{$\displaystyle #2$}}
	{\colorbox{#1}{$\textstyle #2$}}
	{}
	{}}

\newcounter{probcount}% new counter for problem numbers, starts at 0 by default
\NewDocumentEnvironment{ problem } {O{} +b}% Problem environment for homework, autocounts numbers, behaves like a section and can get listed (and hyperlinked) in the toc. 
	{\addtocounter{probcount}{1}% increase counter at the beginning of every problem
	 \phantomsection% invisible page marker for hyperref
	 \addcontentsline{toc}{subsection}{Problem \theprobcount~{\it#1}}% add "Problem \theprobcount. \textit{#1}" to table of contents (toc) as if it were a subsection
	 \begin{trivlist}% begin unmarked ("trivial") list
	 \item {\bf Problem \theprobcount} {\it#1}% print Problem with the counter and optional text 
	 \item #2% +b is for inputs that are long
	 \end{trivlist}% end list
	}{}% don't do anything, then "problem" environment ends, irrelevant because trivlist is already ended
	
\def\UM{\colorbox{blue}{\textcolor{Gold1}{\textbf{University of Michigan}}}} % School Pride
\def\EC{\href{https://github.com/CarpenterEvan/PhysicsReview}{\(\mathbb{E}\textsc{van}~\mathbb{C}\textsc{arpenter}\)}}
\endinput